\section{Circular-phase trace schemes and the remaining signed boundary}
\label{sec:phase-clusters-20261019}

The 16--19 October continuation concerns auxiliary phase ensembles, not
the fixed-coupling BGS conclusion.  The maintained sources and exact proof
revisions are indexed in
\path{mathematics/phase_ensemble_continuation_20261019.md} and
\path{record/chat_integration_20261019.md}.  The new statements have author
proofs and self-review, with exact finite computations where specified;
independent review and priority are not established.

For an odd prime $p$, a nondegenerate alternating form $\Omega$ and
$N=p^{2d}$, put
\[
 S_{xy}=N^{-1/2}\exp(2\pi i\Omega(x,y)/p),\qquad U=DS,
 \qquad V_m=\mathbb E|\operatorname{Tr}U^m|^2,
\]
where $D$ has independent uniform phases.  Then $S=S^*$ and $S^2=I$.
For every fixed number of trace powers the joint law is
\[
 (\operatorname{Tr}U,\operatorname{Tr}U^2,\ldots,
       \operatorname{Tr}U^M)
 \ \Longrightarrow\ (Z,Z^2,\sqrt3Z_3,\ldots,\sqrt M Z_M).
\]
The Gaussian variables displayed separately are independent.  The exact
identity $\operatorname{Tr}U^2=(\operatorname{Tr}U)^2$ prevents identifying
the first two global traces with independent CUE traces.  This does not
by itself decide the local microscopic law.  The polynomial trace-window
results likewise do not reach microscopic absorption.

\subsection{Statements of TRACE-095--097 and TRACE-099}

The following statements (canonical source:
\path{mathematics/phase_ensemble_continuation_20261019.md}) are recorded
here because this section is their registered manuscript source.

TRACE-095 (conditional). Under the pinned feedback trajectory, covariance
and moment inputs and the reversal hypothesis of the source, the Gaussian
limit of every fixed selected block extends to all $u_N\to\infty$ with
$u_N/N\to0$, with error $O((u_N/N)^{1/4}+u_N^{-1/8})$ for the stated
smooth bounded tests of disjoint entries. The substitution $u=2$ is not
permitted.

TRACE-096. For a deterministic unitary $S$ with independent row phases,
spectral averaging identifies the law of the open-port vector $X_i$ with
the eigenvector-ratio law weighted by $|\psi_l(i)|^2$. Closing one port
of a two-port cavity with transmission $T$ by a fresh uniform phase gives
$J=T/|1-\sqrt{1-T}\,e^{i\phi}|^2$ and $Y=J/(1+J)$; the resulting
probability kernel is weakly continuous and injective, so the transmission
laws converge to $\mathrm{Uniform}[0,1]$ if and only if the corresponding
$Y$ laws converge to $\mathrm{Beta}(1,2)$. This is an equivalence
criterion; it does not prove either convergence.

TRACE-097 (conditional). Conditional on the fixed-level TV compactness
input TC of the source (the full $G$ or rank condition $R$ stated there),
the complete normalized positive path-clock current is controlled at
arbitrary cutoffs $K_N\to\infty$. The statement evaluates neither the
finite-channel transmission law nor the complete signed trace sum, and
TC is not discharged in this section.

TRACE-099. For the symplectic seeds the signed mixed-moment comparison
reaches trace powers $o(N^{1/5})$; for generic fixed bilinear seeds of
dimension $b$ it reaches $o(N^{1/(6b)})$, with the actual Gauss phases
retained. The associated disk laws concern fixed compact subsets of
$|z|<1$; they do not reach $z=e^{-s/N}$ or microscopic absorption.

\subsection{Exact schemes and a joint quartic limit}

The signed phase expansion assigns order
$j=c+\tfrac12\operatorname{rank}(H_\Pi^{\mathsf T}KH_\Pi)$ and defect
$\delta=2j-e$ to an open connected cluster with $e$ common-edge ports.
Connected clusters satisfy $j\le\delta+1$.  Exact neutral peeling leaves
primitive order at most $2\delta$ and size at most $16\delta$.  This
controls each fixed defect without summing all connected types.

The saturated-support classification through rank six bounds the connected
four-port core by size sixteen.  Exact enumeration gives a factor $1-z$
in each of its twenty-four open channels.  The physical quotient at
$z=1$ is $-191/12$.  Together with the earlier primitives this completes
the total-defect-two sector for $0<\tau<1$, where $m/N\to\tau$.
For fixed $h\ge3$, the stable coefficient has degree at most $2h-3$
when $m>8h$; its vanishing is not asserted.

There is also a different joint family: all neutral additions and arbitrary
numbers of connected clusters of individual defects one or two.  Its
correction $\mathcal R_{\le2}$ is normalized to $V_m/m$. The gluing uses
Stanley~\cite{Stanley1981Cycles}, Lemmas~2.1--2.2 and Theorem~3.1;
the project supplies the port weights, cumulants and uniform majorant. Define
\[
 v(x)=\frac7{12}+\frac53x,\quad
 q_3(x)=2+\frac32x+\frac12x^2,\quad
 q_4(x)=6+\frac{25}2x+\frac{19}2x^2+2x^3.
\]
For each $x\in[0,1]$, let $w_x$ be the unique positive solution of
\[
 \tau=w_x+w_x^2+\frac{q_3(x)}2w_x^3+\frac{q_4(x)}6w_x^4,
 \qquad D_x=1+2w_x+\frac32q_3(x)w_x^2+\frac23q_4(x)w_x^3.
\]
For $0<\tau<1/9$ the exact five-index sum has limit
\[
 L_{\le2}(\tau)=\int_0^1
 \left[\exp\{w_x+v(x)w_x^2\}\frac{\tau}{w_xD_x}-1\right]dx.
\]
One positive majorant controls all five indices, with tail
$11(9\tau_0)^{H+1}$ on $m/N\le\tau_0<1/9$ and a uniform
$O_{\tau_0}(N^{-1})$ finite-size error for $9\le m\le\tau_0N$.
The result includes the earlier overlapping cluster families and must not
be added to them again.  Its expansion begins
\[
 L_{\le2}(\tau)=-3\tau^3+\frac{20893}{864}\tau^4+O(\tau^5).
\]
The omitted individual defects remain essential.  The connected
$(j,\delta,e)=(3,3,3)$ source is known only through a finite size-nine
prefix in the latest check.  Its full value is not assigned by fitting
the desired CUE answer.

\subsection{What exact Cauchy averaging does not identify}

For any Hermitian involution the local iid-phase problem admits the
quantitative Cayley comparison with $S-\operatorname{diag}(\mathrm{Cauchy})$
stated in the maintained source.  Universality of this particular dense
deterministic hopping model is still required.

An exact finite-dimensional consequence of bounded holomorphic averaging
is useful for separating that requirement.  If $B=B^*$, $E$ is real,
$T$ is iid standard Cauchy diagonal and $A\ge0$ is nonzero, then
\[
 Q_A=\operatorname{Tr}\bigl(A(B-T-EI)^{-1}\bigr)
\]
is Cauchy with complex parameter
$\operatorname{Tr}(A(B-EI+iI)^{-1})$ in the lower half-plane.
For $B=S=S^{-1}$ and $E=0$, its centre and scale are
$\operatorname{Tr}(AS)/2$ and $\operatorname{Tr}(A)/2$.
The contour is applied to a bounded transform of $Q_A$, not to its
undefined first moment.  The general Cauchy/Herglotz principle is
classical; see Aizenman--Warzel,
\href{https://arxiv.org/abs/1312.7769}{arXiv:1312.7769v2}.

Even all these positive projections do not determine a matrix law.  For
the two-dimensional exchange involution, compare
\[
 X=(S-\operatorname{diag}(t_1,t_2))^{-1},\qquad
 Y=S/2+\operatorname{diag}(c_1,c_2)/2,
\]
with independent standard Cauchy variables.  Every positive projection
$\operatorname{Tr}(AX)$ and $\operatorname{Tr}(AY)$ agrees in law.
For $A=\operatorname{diag}(1,-1)$ the signed tail constants are instead
$4/\pi^2$ and $2/\pi$.  More generally, for every non-diagonal Hermitian involution, including
the exact Fourier seed, $H^{-1}$ and $H/2$ have identical positive scalar
projections but matrix laws at total variation distance one. A rank-two
signed source separates a deterministic value from a probability-zero
polynomial level set. This obstructs the proposed positive-source closure,
not BGS. Finite-dimensional total-variation separation does not exclude
coincident limits for particular spectral observables.

All higher connected types, the complete signed mixed hierarchy and the
comparison with the unchanged frozen propagator remain open.  The
conditional feedback and TV-compactness routes retain their named inputs.
